% Unit 042 terjemahan Bahasa Indonesia dari chapter3.tex baris 1882--2214.
% Sumber: Methods of Algebra, Volume 2, commit
% 9a5803ff77dd3257484cb177f851a73770a59dd3 (CC BY 4.0).
% stable-unit-id: o014.aljabr2.chapter3.resolutions
% Cakupan: Bagian 3.11 lengkap.

% segment-id: o014.li2.chapter3.resolutions.h001
\section{Resolusi}\label{sec:resolutions}
\hypertarget{o014.aljabr2.chapter3.resolutions}{ }
% segment-id: o014.li2.chapter3.resolutions.p001
Misalkan $\mathcal{A}$ suatu kategori abelian. Tinjau kompleks $X \in \Obj(\cate{C}(\mathcal{A}))$. Jika $X^n$ merupakan objek injektif (atau objek projektif) di $\mathcal{A}$ untuk setiap $n \in \Z$, maka $X$ dikatakan terdiri atas objek-objek injektif (atau objek-objek projektif). Jangan lupakan fakta sepele bahwa $0$ sekaligus merupakan objek injektif dan objek projektif.

% segment-id: o014.li2.chapter3.resolutions.de001
\begin{definisi}\label{def:resolutions}
	\index{resolusi@resolusi (resolution)}
	\index{resolusi!injektif (injective)}
	\index{resolusi!projektif (projective)}
	Misalkan $X \in \Obj(\cate{C}(\mathcal{A}))$.
% segment-id: o014.li2.chapter3.resolutions.l001
	\begin{itemize}
% segment-id: o014.li2.chapter3.resolutions.i001
		\item Misalkan $X \to I$ suatu kuasi-isomorfisme di $\cate{C}(\mathcal{A})$, dengan $I \in \Obj(\cate{C}^+(\mathcal{A}))$ terdiri atas objek-objek injektif. Maka $X \to I$ disebut \emph{resolusi injektif} dari $X$.
% segment-id: o014.li2.chapter3.resolutions.i002
		\item Misalkan $P \to X$ suatu kuasi-isomorfisme di $\cate{C}(\mathcal{A})$, dengan $P \in \Obj(\cate{C}^-(\mathcal{A}))$ terdiri atas objek-objek projektif. Maka $P \to X$ disebut \emph{resolusi projektif} dari $X$.
	\end{itemize}
\end{definisi}

% segment-id: o014.li2.chapter3.resolutions.p002
Dengan mengganti $\mathcal{A}$ oleh $\mathcal{A}^{\opp}$, yakni dengan membalik arah panah, tampak bahwa resolusi injektif dan resolusi projektif merupakan konsep yang saling dual.

% segment-id: o014.li2.chapter3.resolutions.ex001
\begin{contoh}\label{eg:resolution-single}
	Sebagai kasus khusus yang paling mendasar sekaligus paling klasik, tinjau resolusi injektif $X \to I$ dari $X \in \Obj(\mathcal{A})$. Definisi hanya mensyaratkan $I \in \Obj(\cate{C}^+(\mathcal{A}))$, tetapi kita ingin mengambil $I \in \Obj(\cate{C}^{\geq 0}(\mathcal{A}))$. Hal ini dapat dilakukan dengan menggunakan pemenggalan
% segment-id: o014.li2.chapter3.resolutions.q001
	\[ X \to I \to \tau^{\geq 0} I . \]
	Untuk menunjukkan bahwa komposisinya merupakan resolusi injektif, dua sifat harus diperiksa. Pertama, Lema \ref{prop:truncation-ses} menunjukkan bahwa komposisi ini tetap merupakan kuasi-isomorfisme. Kedua, $\tau^{\geq 0} I$ tetap terdiri atas objek-objek injektif. Mengapa? Tinjau barisan eksak pendek $0 \to \Image\left(d_I^{n-1} \right) \to I^n \to \Coker\left( d_I^{n-1} \right) \to 0$. Karena $d_I^{n-1} = 0$ untuk $n \ll 0$, sedangkan $\Coker\left( d_I^{n-1} \right) \simeq \Image\left( d_I^n \right)$ untuk $n < 0$, Lema \ref{prop:inj-proj-split} dan \ref{prop:prod-injective-projective} secara rekursif menunjukkan bahwa untuk $n \leq 0$,
% segment-id: o014.li2.chapter3.resolutions.q002
	\[ I^n \simeq \Image\left(d_I^{n-1} \right) \oplus \Coker\left( d_I^{n-1} \right) , \quad \text{sehingga}\; \Coker\left(d_I^{n-1} \right) \;\text{merupakan objek injektif}; \]
	titik awal rekursi adalah $n \ll 0$, ketika rumus di atas menjadi $0 = 0 \oplus 0$. Dengan mengambil $n=0$, diperoleh bahwa $\Coker\left( d_I^{-1}\right)$ merupakan objek injektif. Jadi, $\tau^{\geq 0} I$ terdiri atas objek-objek injektif.

	Karena dalam praktik yang diperhatikan ialah resolusi yang cukup ``dalam'', pengamatan di atas menunjukkan bahwa kita boleh mengganti $X \to I$ dengan $X \to \tau^{\geq 0} I$, sehingga resolusi injektif berbentuk
% segment-id: o014.li2.chapter3.resolutions.g001
	\[\begin{tikzcd}[row sep=small]
		\cdots 0 \arrow[r] & 0 \arrow[r] \arrow[d] & X \arrow[d] \arrow[r] & 0 \arrow[r] \arrow[d] & 0 \arrow[r] \arrow[d] & \cdots \\
		\cdots 0 \arrow[r] & 0 \arrow[r] & I^0 \arrow[r] & I^1 \arrow[r] & I^2 \arrow[r] & \cdots
	\end{tikzcd}\]
	Diagram ini juga dapat diratakan dan dipandang sebagai barisan eksak di $\mathcal{A}$,
% segment-id: o014.li2.chapter3.resolutions.q003
	\[ 0 \to X \to I^0 \to I^1 \to \cdots, \quad \text{setiap $I^n$ merupakan objek injektif, $n \geq 0$}. \]
	
	Demikian pula, untuk resolusi projektif $P \to X$, dengan mengganti $P \to X$ oleh komposisi $\tau^{\leq 0} P \to P \to X$, kita boleh menganggapnya berasal dari barisan eksak di $\mathcal{A}$,
% segment-id: o014.li2.chapter3.resolutions.q004
	\[ \cdots \to P^{-1} \to P^0 \to X \to 0, \quad \text{setiap $P^n$ merupakan objek projektif, $n \leq 0$.} \]
	Rumus di atas lazim ditulis dalam notasi kompleks rantai sebagai $\cdots \to P_1 \to P_0 \to X \to 0$, dengan $P_n := P^{-n}$.
\end{contoh}

% segment-id: o014.li2.chapter3.resolutions.p003
Persoalannya ialah keberadaan resolusi injektif atau projektif. Berdasarkan dualitas, kita cukup terlebih dahulu membahas konstruksi resolusi injektif. Kasus yang paling sederhana tetaplah $X \in \Obj(\mathcal{A})$. Andaikan $\mathcal{A}$ memiliki cukup banyak objek injektif; lihat Definisi \ref{def:enough-injectives}. Pertama-tama, pilih sembarang monomorfisme $X \to I^0$ dengan $I^0$ objek injektif. Secara rekursif, andaikan telah diperoleh barisan eksak
% segment-id: o014.li2.chapter3.resolutions.q005
\[ 0 \to X \to I^0 \to \cdots \to I^n, \quad n \geq 0, \]
dengan $I^0, \ldots, I^n$ semuanya objek injektif. Pada langkah awal $n=0$, objek kokernel yang digunakan ialah $\Coker\left[ X \to I^0\right]$; untuk $n \geq 1$, objek itu ialah $\Coker\left[ I^{n-1} \to I^n\right]$. Dalam kedua kasus, pilih objek injektif $I^{n+1}$ dan monomorfisme dari objek kokernel tersebut ke $I^{n+1}$. Dengan mengambil komposisinya, kita memperoleh $I^n \to I^{n+1}$ sedemikian sehingga $0 \to X \to \cdots \to I^{n+1}$ juga eksak. Mengulangi operasi ini menghasilkan resolusi injektif dari $X$.
\footnote{Catatan penerjemah (O014-C056): sumber menerapkan rumus kokernel
$\Coker[I^{n-1}\to I^n]$ juga pada $n=0$, sehingga memunculkan $I^{-1}$
yang belum terdefinisi. Target memakai $\Coker[X\to I^0]$ untuk langkah awal
dan membatasi rumus tercetak pada $n\geq 1$.}

% segment-id: o014.li2.chapter3.resolutions.p004
Dengan asumsi bahwa $\mathcal{A}$ memiliki cukup banyak objek projektif, konstruksi resolusi projektif $\cdots \to P^0 \to X \to 0$ dari $X \in \Obj(\mathcal{A})$ sepenuhnya dual. Konstruksi-konstruksi ini dapat diperluas ke kompleks, tetapi memerlukan beberapa syarat keterbatasan, sebagaimana dirinci berikut ini.

% segment-id: o014.li2.chapter3.resolutions.th001
\begin{teorema}\label{prop:resolution-existence}
	Misalkan $\mathcal{A}^\flat$ suatu subkategori penuh aditif dari $\mathcal{A}$, dan untuk setiap objek $A$ di $\mathcal{A}$ terdapat monomorfisme $A \hookrightarrow B$ (atau epimorfisme $B \twoheadrightarrow A$) sedemikian sehingga $B \in \Obj(\mathcal{A}^\flat)$. Maka,
% segment-id: o014.li2.chapter3.resolutions.l002
	\begin{enumerate}[(i)]
% segment-id: o014.li2.chapter3.resolutions.i003
		\item untuk setiap objek $X$ di $\cate{C}^+(\mathcal{A})$ (atau $\cate{C}^-(\mathcal{A})$), terdapat kuasi-isomorfisme
% segment-id: o014.li2.chapter3.resolutions.q006
		\[ f: X \to I \quad \text{(atau $f: P \to X$)}, \]
		sedemikian sehingga $f$ merupakan monomorfisme (atau epimorfisme), dan $I \in \Obj(\cate{C}^+(\mathcal{A}^\flat))$ (atau $P \in \Obj(\cate{C}^-(\mathcal{A}^\flat))$).
% segment-id: o014.li2.chapter3.resolutions.i004
		\item Lebih tepatnya, jika $m \in \Z$ dan $X \in \Obj(\cate{C}^{\geq m}(\mathcal{A}))$ (atau $X \in \Obj(\cate{C}^{\leq m}(\mathcal{A}))$), maka dapat dipilih $I \in \Obj(\cate{C}^{\geq m}(\mathcal{A}))$ (atau $P \in \Obj(\cate{C}^{\leq m}(\mathcal{A}))$).
	\end{enumerate}

	Sebagai akibat, jika $\mathcal{A}$ memiliki cukup banyak objek injektif (atau objek projektif), maka $X$ seperti di atas memiliki resolusi injektif $X \hookrightarrow I$ (atau resolusi projektif $P \twoheadrightarrow X$).
\end{teorema}
% segment-id: o014.li2.chapter3.resolutions.pf001
\begin{bukti}
	Berdasarkan dualitas, berikut ini hanya dibahas versi dengan $X$ berasal dari $\cate{C}^+(\mathcal{A})$. Untuk (i), secara rekursif andaikan telah diperoleh diagram komutatif
% segment-id: o014.li2.chapter3.resolutions.g002
	\[\begin{tikzcd}
		\cdots \arrow[r] & X^{n-1} \arrow[r, "{d_X^{n-1}}"] \arrow[d, "{f^{n-1}}"'] & X^n \arrow[r, "{d_X^n}"] \arrow[d, "{f^n}"] & X^{n+1} \arrow[r] & \cdots \\
		\cdots \arrow[r] & I^{n-1} \arrow[r, "{d_I^{n-1}}"'] & I^n & &
	\end{tikzcd}\]
	sedemikian sehingga baris kedua merupakan kompleks yang terdiri atas objek-objek di $\mathcal{A}^\flat$, $m \ll 0 \implies I^m = 0$, setiap $f^m$ merupakan monomorfisme ($m \leq n$), dan untuk $m \leq n-1$, $f^m$ menginduksi isomorfisme $\Ker\left( d_X^m \right) / \Image\left( d_X^{m-1} \right) \rightiso \Ker\left( d_I^m \right) / \Image\left( d_I^{m-1} \right)$.

	Kita ingin memperpanjang baris kedua diagram ke kanan. Pertama-tama andaikan
% segment-id: o014.li2.chapter3.resolutions.q007
	\begin{equation}\label{eqn:resolution-existence-aux0}
		f^n\left( \Image\left(d_X^{n-1} \right) \right) = \Image\left( d^{n-1}_I \right) \cap f^n\left( \Ker\left( d_X^n \right) \right) = \Image\left( d^{n-1}_I \right) \cap f^n(X^n) ;
	\end{equation}
	perhatikan bahwa relasi inklusi $\cdots \subset \cdots \subset \cdots$ di sini selalu berlaku. Konstruksi dorong keluar memberikan diagram
% segment-id: o014.li2.chapter3.resolutions.g003
	\begin{equation*}\begin{tikzcd}
		X^n / \Ker\left( d_X^n \right) \arrow[hookrightarrow, r, "\delta: \text{diinduksi oleh $d_X^n$}"] \arrow[d, "\alpha: \text{diinduksi oleh $f^n$}"'] & X^{n+1} \arrow[d, "\beta"] \\
		I^n / \left( f^n \left( \Ker\left( d_X^n \right) \right) + \Image\left( d_I^{n-1}\right) \right) \arrow[r, "\eta"'] & T \arrow[phantom, lu, "\boxplus" description]
	\end{tikzcd}\end{equation*}
	Perhatikan bahwa $f^n$ merupakan monomorfisme, sedangkan \eqref{eqn:resolution-existence-aux0} mengakibatkan
% segment-id: o014.li2.chapter3.resolutions.q008
	\begin{multline*}
		\left( f^n\left( \Ker\left( d_X^n \right) \right) + \Image\left( d_I^{n-1}\right) \right) \cap f^n(X^n) \\
		\xlongequal{\text{Teorema \ref{prop:subobject-modularity}}} f^n\left( \Ker\left( d_X^n \right) \right) + \left( \Image\left( d_I^{n-1}\right) \cap f^n(X^n) \right) \\
		= f^n\left( \Ker\left( d_X^n \right) \right) + f^n\left(\Image\left(d_X^{n-1} \right) \right) = f^n\left( \Ker\left( d_X^n \right) \right),
	\end{multline*}
	sehingga operasi-operasi dasar yang dibahas dalam \S\ref{sec:Abel-cat-subobjects} menunjukkan bahwa $\alpha$ juga merupakan monomorfisme.
	
	Dengan menerapkan Proposisi \ref{prop:Abel-cat-pull-push} dua kali, tampak bahwa $\beta$ dan $\eta$ dalam diagram dorong keluar juga merupakan monomorfisme. Pilih monomorfisme $T \hookrightarrow I^{n+1}$ sedemikian sehingga $I^{n+1} \in \Obj(\mathcal{A}^\flat)$. Darinya diperoleh morfisme komposit
% segment-id: o014.li2.chapter3.resolutions.q009
	\begin{gather*}
		d_I^n: I^n \twoheadrightarrow I^n / \left( f^n \left( \Ker\left( d_X^n \right) \right) + \Image\left( d_I^{n-1}\right) \right) \xrightarrow{\eta} T \hookrightarrow I^{n+1}, \\
		f^{n+1}: X^{n+1} \xrightarrow{\beta} T \hookrightarrow I^{n+1}.
	\end{gather*}
	Jelas bahwa $d_I^n d_I^{n-1} = 0$, bahwa $f^{n+1}$ merupakan monomorfisme, dan bahwa $d_I^n f^n = f^{n+1} d_X^n$. Monomorfisme $f^n$ menginduksi pada kohomologi
% segment-id: o014.li2.chapter3.resolutions.q010
	\begin{align*}
		\frac{\Ker\left( d_X^n \right)}{\Image\left( d_X^{n-1} \right)} \longrightarrow \frac{\Ker\left( d_I^n \right)}{\Image\left( d_I^{n-1} \right)} & =  \frac{ f^n\left(\Ker d_X^n \right) + \Image\left( d_I^{n-1} \right)}{ \Image\left(d_I^{n-1}\right) } \\
		\text{($\because$ Teorema \ref{prop:Abel-cat-isom-thm} (iii))} \quad & \simeq \frac{ f^n\left(\Ker d_X^n \right) }{ f^n\left(\Ker d_X^n \right) \cap \Image\left(d_I^{n-1}\right) };
	\end{align*}
	berdasarkan \eqref{eqn:resolution-existence-aux0} dan sifat monomorfisme $f^n$, pemetaan ini jelas merupakan isomorfisme. Perpanjangan ke kanan berhasil.

	Sekarang akan dijelaskan cara memodifikasi $I^n$ dalam keadaan umum agar tereduksi ke kasus ketika \eqref{eqn:resolution-existence-aux0} berlaku. Tinjau morfisme kanonik
% segment-id: o014.li2.chapter3.resolutions.g004
	\[\begin{tikzcd}
		I^n \arrow[hookrightarrow, shift left, r, "\iota_1"] & I^n \oplus \left( X^n / \Image\left( d_X^{n-1} \right) \right) \arrow[twoheadrightarrow, shift left, l, "p_1"] \arrow[twoheadrightarrow, shift right, r, "p_2"'] & X^n / \Image\left( d_X^{n-1} \right) \arrow[hookrightarrow, shift right, l, "\iota_2"'] & X^n \arrow[twoheadrightarrow, l, "q"'] .
	\end{tikzcd}\]
	Pilih sembarang $J \in \Obj(\mathcal{A}^\flat)$ dan monomorfisme $j: I^n \oplus \left( X^n / \Image\left( d_X^{n-1} \right) \right) \hookrightarrow J$. Definisikan morfisme komposit
% segment-id: o014.li2.chapter3.resolutions.q011
	\begin{gather*}
		f' : X^n \xrightarrow{(f^n, q)} I^n \oplus \left( X^n / \Image\left( d_X^{n-1} \right) \right) \xrightarrow{j} J, \\
		d': I^{n-1} \xrightarrow{d_I^{n-1}} I^n \xrightarrow{\iota_1} I^n \oplus \left( X^n / \Image\left( d_X^{n-1} \right) \right) \xrightarrow{j} J.
	\end{gather*}
	Mudah dilihat bahwa $f'$ merupakan monomorfisme, $\Ker(d') = \Ker\left( d_I^{n-1} \right)$, dan $f' d_X^{n-1} = d' f^{n-1}$.
	
	Selain itu, dengan bekerja di $I^n \oplus \left( X^n / \Image\left( d_X^{n-1} \right) \right)$, tampak bahwa
% segment-id: o014.li2.chapter3.resolutions.q012
	\[ f'(X^n) \cap \Image(d') = f'\left( \Image\left( d_X^{n-1} \right) \right). \]
	Hal ini menunjukkan bahwa dengan mengganti $d_I^{n-1}: I^{n-1} \to I^n$ (atau $f^n: X^n \to I^n$) oleh $d': I^{n-1} \to J$ (atau $f': X^n \to J$), kita dapat memastikan bahwa \eqref{eqn:resolution-existence-aux0} berlaku.
	
	Terakhir, kita membahas (ii). Dalam konstruksi di atas, ambil $\cdots = I^{m-2} = I^{m-1} = 0$. Jelas bahwa \eqref{eqn:resolution-existence-aux0} berlaku untuk $n = m-1$ (semua sukunya nol), sehingga proses memperpanjang diagram ke kanan tidak menghasilkan suku tak nol berderajat $< m$. Ini membuktikan pernyataan yang diinginkan.
\end{bukti}

% segment-id: o014.li2.chapter3.resolutions.co001
\begin{korolari}
	Misalkan $\mathcal{A}$ memiliki cukup banyak objek injektif (atau objek projektif). Maka, $X \in \Obj(\cate{C}(\mathcal{A}))$ memiliki resolusi injektif $X \to I$ (atau resolusi projektif $P \to X$) jika dan hanya jika $\Hm^n(X) = 0$ untuk $n \ll 0$ (atau $n \gg 0$).
\end{korolari}
% segment-id: o014.li2.chapter3.resolutions.pf002
\begin{bukti}
	Ambil kasus resolusi injektif sebagai contoh. Syarat kuasi-isomorfisme dalam definisi resolusi memberikan arah ``hanya jika''. Di sisi lain, untuk $m \ll 0$, funktor pemenggalan memberikan kuasi-isomorfisme $X \to \tau^{\geq m} X \in \Obj(\cate{C}^+(\mathcal{A}))$. Dengan demikian, arah ``jika'' tereduksi ke Proposisi \ref{prop:resolution-existence}.
\end{bukti}

% segment-id: o014.li2.chapter3.resolutions.p005
Nilai teoretis resolusi injektif atau projektif dapat dilihat dari hasil dasar berikut, yang juga mengakibatkan bahwa resolusi injektif atau projektif unik hingga homotopi.

% segment-id: o014.li2.chapter3.resolutions.th002
\begin{teorema}\label{prop:resolution-homotopy}
	Tetapkan suatu kuasi-isomorfisme $\alpha: X \to Y$ di $\cate{C}(\mathcal{A})$.
% segment-id: o014.li2.chapter3.resolutions.l003
	\begin{itemize}
% segment-id: o014.li2.chapter3.resolutions.i005
		\item Diberikan morfisme $\gamma: X \to I$, dengan $I$ suatu objek di $\cate{C}^+(\mathcal{A})$ yang terdiri atas objek-objek injektif. Di $\cate{K}(\mathcal{A})$, terdapat tepat satu morfisme $\beta$ yang membuat diagram berikut komutatif:
% segment-id: o014.li2.chapter3.resolutions.g005
		\[\begin{tikzcd}[row sep=tiny]
			& Y \arrow[dd, "\beta"] \\
			X \arrow[ru, "\alpha"] \arrow[rd, "\gamma"'] & \\
			& I
		\end{tikzcd}\]
% segment-id: o014.li2.chapter3.resolutions.i006
		\item Diberikan morfisme $\gamma: P \to Y$, dengan $P$ suatu objek di $\cate{C}^-(\mathcal{A})$ yang terdiri atas objek-objek projektif. Di $\cate{K}(\mathcal{A})$, terdapat tepat satu morfisme $\beta$ yang membuat diagram berikut komutatif:
% segment-id: o014.li2.chapter3.resolutions.g006
		\[\begin{tikzcd}[row sep=tiny]
			& X \arrow[ld, "\alpha"'] \\
			Y & \\
			& P \arrow[lu, "\gamma"] \arrow[uu, "\beta"']
		\end{tikzcd}\]
	\end{itemize}
\end{teorema}

% segment-id: o014.li2.chapter3.resolutions.p006
% future-source-xref: prop:cplx-triangulated (Teorema 4.4.1)
Kedua pernyataan tersebut tentu saling dual. Setelah Teorema 4.4.1 dalam sumber (bagian yang belum tercakup dalam pembaca ini), kelak akan diberikan pembahasan yang ringkas. Jika $\alpha$ merupakan monomorfisme (atau epimorfisme), kita bahkan dapat memilih $\beta$ lebih lanjut sehingga diagram pertama (atau kedua) komutatif di $\cate{C}(\mathcal{A})$; ini adalah isi Lema \ref{prop:ext-alpha-beta} berikut beserta versi dualnya. Sesungguhnya, bagian latihan bab ini akan memberikan pembuktian langsung atas Teorema \ref{prop:resolution-homotopy} berdasarkan lema tersebut, lengkap dengan petunjuk; pembaca dianjurkan untuk mencobanya.

% segment-id: o014.li2.chapter3.resolutions.p007
Pendekatan mana pun yang dipilih tidak dapat lepas dari sifat berikut.

% segment-id: o014.li2.chapter3.resolutions.le001
\begin{lema}\label{prop:resolution-homotopy-aux}
	Misalkan objek $X$ di $\cate{C}(\mathcal{A})$ asiklik.
% segment-id: o014.li2.chapter3.resolutions.l004
	\begin{itemize}
% segment-id: o014.li2.chapter3.resolutions.i007
		\item Jika objek $I$ di $\cate{C}^+(\mathcal{A})$ terdiri atas objek-objek injektif, maka $\Hom_{\cate{K}(\mathcal{A})}(X, I) = 0$.
% segment-id: o014.li2.chapter3.resolutions.i008
		\item Jika objek $P$ di $\cate{C}^-(\mathcal{A})$ terdiri atas objek-objek projektif, maka $\Hom_{\cate{K}(\mathcal{A})}(P, X) = 0$.
	\end{itemize}
\end{lema}
% segment-id: o014.li2.chapter3.resolutions.pf003
\begin{bukti}
	Cukuplah menangani kasus $I$. Tetapkan $\alpha \in \Hom_{\cate{C}(\mathcal{A})}(X,I)$. Kita mencari suatu keluarga morfisme $h^n: X^n \to I^{n-1}$ sedemikian sehingga
% segment-id: o014.li2.chapter3.resolutions.q013
	\begin{equation}\label{eqn:alpha-homotopy-aux}
		\alpha^n = d_I^{n-1} h^n + h^{n+1} d_X^n , \quad n \in \Z.
	\end{equation}

	Secara rekursif, andaikan telah diperoleh $\ldots, h^{n-1}, h^n$ sedemikian sehingga \eqref{eqn:alpha-homotopy-aux} berlaku untuk $\ldots, \alpha^{n-2}, \alpha^{n-1}$. Andaian ini beralasan karena kita dapat mengambil $\ldots, h^{n-1}, h^n$ semuanya $0$ ketika $n \ll 0$. Kita mencari morfisme yang ditandai dengan garis putus-putus,
% segment-id: o014.li2.chapter3.resolutions.g007
	\[\begin{tikzcd}[row sep=large]
		& I^n & \\
		X^n \arrow[twoheadrightarrow, r, "{d_X^n}"'] \arrow[ru, "{\alpha^n - d_I^{n-1} h^n}"] & \Image\left( d_X^n \right) \arrow[dashed, u, "\beta"] \arrow[hookrightarrow, r] & X^{n+1} \arrow[dashed, lu, "{h^{n+1}}"']
	\end{tikzcd}\]
\footnote{Catatan penerjemah (O014-C057): simpul kiri bawah pada sumber
tercetak sebagai $X$, padahal $d_X^n$, $\alpha^n$, dan $h^n$ semuanya
berdomain $X^n$. Target memulihkan simpul tersebut menjadi $X^n$.}
	sedemikian sehingga seluruh diagram komutatif. Pertama,
% segment-id: o014.li2.chapter3.resolutions.q014
	\[ \left( \alpha^n - d_I^{n-1} h^n \right) d_X^{n-1} = d_I^{n-1} \alpha^{n-1} - d_I^{n-1} \left( \alpha^{n-1} - d_I^{n-2} h^{n-1} \right) = 0. \]
	Karena $X$ asiklik, hal ini menginduksi morfisme $\beta$ pada diagram. Selanjutnya, karena $I^n$ merupakan objek injektif, $\beta$ dapat diperluas menjadi $h^{n+1}: X^{n+1} \to I^n$. Ini membuktikan pernyataan yang diinginkan.
\end{bukti}

% segment-id: o014.li2.chapter3.resolutions.p008
% future-source-xref: sec:double-cplx-ss (Bagian 5.6)
Dari sudut pandang kategori turunan, informasi yang diberikan oleh Teorema \ref{prop:resolution-existence} dan \ref{prop:resolution-homotopy} sudah memadai untuk mulai mempelajari funktor turunan. Beberapa jenis resolusi yang akan segera diperkenalkan lebih sesuai digunakan bersama barisan spektral; lihat Bagian 5.6 dalam sumber (belum tercakup dalam pembaca ini). Berikut ini kita merumuskan versi injektif. Kita memerlukan beberapa persiapan.

% segment-id: o014.li2.chapter3.resolutions.le002
\begin{lema}\label{prop:ext-alpha-beta}
	Misalkan $\alpha: X \to Y$ suatu kuasi-isomorfisme di $\cate{C}(\mathcal{A})$, dan $\alpha$ merupakan monomorfisme. Diberikan suatu morfisme $\gamma: X \to I$ di $\cate{C}(\mathcal{A})$, dengan $I \in \Obj(\cate{C}^+(\mathcal{A}))$ terdiri atas objek-objek injektif. Terdapat morfisme $\beta: Y \to I$ di $\cate{C}(\mathcal{A})$ sedemikian sehingga $\gamma = \beta\alpha$.
\end{lema}
% segment-id: o014.li2.chapter3.resolutions.pf004
\begin{bukti}
	Dengan bantuan $\alpha^n$, pandang setiap $X^n$ sebagai subobjek dari $Y^n$; selanjutnya $\alpha^n$ dihilangkan dari notasi. Kita membangun $\beta^n$ langkah demi langkah; untuk $n \ll 0$, tetapkan $\beta^n = 0$. Andaikan telah diperoleh $\ldots, \beta^{n-2}, \beta^{n-1}$ yang memenuhi syarat. Morfisme yang dicari, $\beta^n: Y^n \to I^n$, dapat dan harus ditentukan pada subobjek $X^n$ dan $\Image(d^{n-1}_Y)$ sebagai berikut.
% segment-id: o014.li2.chapter3.resolutions.l005
	\begin{itemize}
% segment-id: o014.li2.chapter3.resolutions.i009
		\item Pada subobjek $X^n$, $\beta^n$ sama dengan $\gamma^n$.
% segment-id: o014.li2.chapter3.resolutions.i010
		\item Pada subobjek $\Image(d^{n-1}_Y)$, $\beta^n$ diinduksi oleh komposisi $Y^{n-1} \xrightarrow{\beta^{n-1}} I^{n-1} \xrightarrow{d_I^{n-1}} I^n$. Agar pernyataan ini bermakna, kita tegaskan bahwa pembatasan $d_I^{n-1} \beta^{n-1}$ pada $\Ker(d_Y^{n-1})$ adalah $0$. Perhatikan diagram komutatif
% segment-id: o014.li2.chapter3.resolutions.g008
		\[\begin{tikzcd}
			\Image(d_X^{n-2}) \arrow[r] \arrow[hookrightarrow, d] & \Ker(d_X^{n-1}) \arrow[r] \arrow[hookrightarrow, d] & X^{n-1} \arrow[rd, "{\gamma^{n-1}}"] \arrow[hookrightarrow, d] & & \\
			\Image(d_Y^{n-2}) \arrow[r] & \Ker(d_Y^{n-1}) \arrow[r] & Y^{n-1} \arrow[r, "{\beta^{n-1}}"'] & I^{n-1} \arrow[r, "{d_I^{n-1}}"'] & I^n . 
		\end{tikzcd}\]
		Dari andaian pada $\ldots, \beta^{n-2}, \beta^{n-1}$ dan $d_I^{n-1} d_I^{n-2} = 0$, diketahui bahwa komposisi pada baris kedua adalah $0$. Selain itu, karena $\gamma$ merupakan morfisme antarkompleks, komposisi $\Ker(d_X^{n-1}) \to \cdots \to I^n$ adalah $0$. Secara keseluruhan, diperoleh diagram komutatif
% segment-id: o014.li2.chapter3.resolutions.g009
		\[\begin{tikzcd}
			\Hm^{n-1}(X) \arrow[d] \arrow[r, "0"] & I^n \\
			\Hm^{n-1}(Y) \arrow[ru, "\text{diinduksi oleh $d_I^{n-1} \beta^{n-1}$}"'] &
		\end{tikzcd}\]
		Namun, syarat kuasi-isomorfisme menunjukkan bahwa $\Hm^{n-1}(X) \rightiso \Hm^{n-1}(Y)$. Penegasan tersebut terbukti.
	\end{itemize}
	
	Jika kita dapat membuktikan kesamaan subobjek
% segment-id: o014.li2.chapter3.resolutions.q015
	\[ X^n \cap \Image\left(d_Y^{n-1} \right) = \Image\left( d_X^{n-1} \right), \]
	maka, berdasarkan diagram dorong keluar dalam Proposisi \ref{prop:biproduct-sum-intersection} beserta sifat universal dorong keluar, morfisme tersebut dapat diperluas ke $X^n + \Image(d^{n-1}_Y)$, lalu dengan menggunakan sifat $I^n$ sebagai objek injektif diperluas lagi menjadi $\beta^n: Y^n \to I^n$. Perhatikan bahwa inklusi $\supset$ sudah jelas.
	
	Pokoknya ialah membuktikan bahwa jika sembarang morfisme $\phi: T \to Y^n$ di $\mathcal{A}$ memfaktor melalui $X^n \cap \Image\left(d_Y^{n-1} \right)$, maka morfisme itu otomatis memfaktor melalui $\Image\left( d_X^{n-1} \right)$. Morfisme $\phi$ semacam itu tentu menginduksi $T \to \Ker\left( d_X^n \right) = X^n \cap \Ker\left( d_Y^n \right)$. Sekarang bentuk komposisi sepanjang baris kedua diagram komutatif
% segment-id: o014.li2.chapter3.resolutions.g010
	\[\begin{tikzcd}
		T \arrow[r, "\phi"] \arrow[d, "\phi"'] & \Ker(d_X^n) \arrow[twoheadrightarrow, r] \arrow[hookrightarrow, d] & \Hm^n(X) \arrow[d] \\
		\Image(d_Y^{n-1}) \arrow[hookrightarrow, r] & \Ker(d_Y^n) \arrow[twoheadrightarrow, r] & \Hm^n(Y)
	\end{tikzcd}\]
	dan hasilnya adalah $0$. Karena $\Hm^n(X) \rightiso \Hm^n(Y)$, komposisi pada baris pertama juga $0$. Dengan demikian, $\phi$ memfaktor melalui $\Image\left( d_X^{n-1} \right)$.
	
	Menurut konstruksi, $\beta = (\beta^n)_n$ memberikan morfisme $Y \to I$, sedangkan syarat $\gamma = \beta\alpha$ juga langsung diperoleh dari konstruksi.
\end{bukti}

% segment-id: o014.li2.chapter3.resolutions.le003
\begin{lema}\label{prop:split-differential}
	Misalkan $A, C \in \Obj(\cate{C}(\mathcal{A}))$. Untuk setiap $n \in \Z$, definisikan $B^n := A^n \oplus C^n$, yang dilengkapi dengan morfisme jelas $A^n \xrightarrow{i^n} B^n \xrightarrow{p^n} C^n$. Kita memiliki bijeksi
% segment-id: o014.li2.chapter3.resolutions.q016
	\[ \begin{gathered}
		\Hom_{\cate{C}(\mathcal{A})}(C, A[1]) \xrightarrow{1:1} \\
		\left\{(d_B^n)_{n \in \Z}\;\middle|\;
		\begin{gathered}
			(B^n, d_B^n)_n \;\text{menjadi kompleks}, \\
			0 \to A \xrightarrow{(i^n)_n} B \xrightarrow{(p^n)_n} C \to 0 \;\text{eksak}
		\end{gathered}\right\}
	\end{gathered} \]
	yang secara eksplisit mengirim $\delta: C \to A[1]$ ke diferensial yang ditentukan dalam notasi matriks sebagai
% segment-id: o014.li2.chapter3.resolutions.q017
	\[ d_B^n = \begin{pmatrix}
		d_A^n & \delta^n \\
		0 & d_C^n
	\end{pmatrix}: B^n \to B^{n+1}. \]
\end{lema}
% segment-id: o014.li2.chapter3.resolutions.pf005
\begin{bukti}
	Setelah diketahui bahwa $(B^n, d_B^n)_n$ merupakan kompleks dan bahwa $(i^n)_n$ serta $(p^n)_n$ merupakan morfisme kompleks, keeksakan $0 \to A \to B \to C \to 0$ dapat diperiksa suku demi suku dan tidak menyisakan persoalan.
	
	Syarat perlu dan cukup agar $d_B^n i^n = i^{n+1} d_A^n$ dan $d_C^n p^n = p^{n+1} d_B^n$ selalu berlaku ialah bahwa terdapat suatu keluarga $\delta^n: C^n \to A^{n+1}$ sedemikian sehingga $d_B^n$ dapat dinyatakan sebagai matriks segitiga atas dalam pernyataan. Persoalannya tereduksi menjadi pencirian $(\delta^n)_n$ yang membuat $d_B^{n+1} d_B^n = 0$. Perhitungan rutin memberikan syarat perlu dan cukup $d_A^{n+1} \delta^n + \delta^{n+1} d_C^n = 0$, yakni $d_{A[1]}^n \delta^n = \delta^{n+1} d_C^n$.
\end{bukti}

% segment-id: o014.li2.chapter3.resolutions.pr001
\begin{proposisi}[Lema Tapal Kuda (Horseshoe Lemma)]\label{prop:horseshoe}
	\index{lematapalkuda@Lema Tapal Kuda (Horseshoe Lemma)}
	Misalkan $0 \to A \to B \to C \to 0$ suatu barisan eksak pendek di $\cate{C}(\mathcal{A})$, dengan $A, B, C \in \Obj(\cate{C}^+(\mathcal{A}))$. Ambil sembarang resolusi injektif $\epsilon: A \to I$ dan $\eta: C \to K$ yang memiliki sifat-sifat seperti dalam Teorema \ref{prop:resolution-existence}. Maka data tersebut dapat diperluas menjadi diagram komutatif yang eksak pada baris-barisnya di $\cate{C}^+(\mathcal{A})$,
% segment-id: o014.li2.chapter3.resolutions.g011
	\[\begin{tikzcd}
		0 \arrow[r] & I \arrow[r] & J \arrow[r] & K \arrow[r] & 0 \\
		0 \arrow[r] & A \arrow[r] \arrow[u, "\epsilon"] & B \arrow[r] \arrow[u, "\kappa"] & C \arrow[r] \arrow[u, "\eta"'] & 0 
	\end{tikzcd}\]
	sedemikian sehingga $\kappa: B \to J$ juga merupakan resolusi injektif, dan $\kappa$ merupakan monomorfisme.
\end{proposisi}

% segment-id: o014.li2.chapter3.resolutions.p009
Pembaca dapat mencoba memberikan pembuktian yang relatif sederhana untuk kasus khusus $A, B, C \in \Obj(\mathcal{A})$.

% segment-id: o014.li2.chapter3.resolutions.pf006
\begin{bukti}
	Bangun diagram dorong keluar di $\cate{C}(\mathcal{A})$,
% segment-id: o014.li2.chapter3.resolutions.g012
	\[\begin{tikzcd}
		0 \arrow[r] & I \arrow[r] & L \arrow[r] \arrow[phantom, ld, "\boxplus" description] & C \arrow[r] & 0 \\
		0 \arrow[r] & A \arrow[r] \arrow[u, "\epsilon"] & B \arrow[r] \arrow[u] & C \arrow[r] \arrow[u, "\identity"'] & 0 
	\end{tikzcd}\]
	Di sini digunakan dua fakta:
% segment-id: o014.li2.chapter3.resolutions.l006
	\begin{compactitem}
% segment-id: o014.li2.chapter3.resolutions.i011
		\item dorong keluar mempertahankan kokernel $C$;
% segment-id: o014.li2.chapter3.resolutions.i012
		\item Proposisi \ref{prop:Abel-cat-pull-push} dan sifat monomorfisme $A \to B$ mengakibatkan bahwa $I \to L$ juga merupakan monomorfisme.
	\end{compactitem}
	Secara khusus, diagram di atas tetap merupakan diagram komutatif yang eksak pada baris-barisnya, sehingga $L \in \Obj(\cate{C}^+(\mathcal{A}))$. Dengan menerapkan lagi Proposisi \ref{prop:Abel-cat-pull-push} dan menggunakan sifat monomorfisme $\epsilon$, diperoleh bahwa $B \to L$ juga merupakan monomorfisme.

	Karena $I^n$ injektif, Lema \ref{prop:inj-proj-split} mengakibatkan bahwa barisan eksak pendek $0 \to I^n \to L^n \to C^n \to 0$ terbelah untuk setiap $n$, sehingga memberikan isomorfisme $\Phi^n: L^n \rightiso I^n \oplus C^n$. Lema \ref{prop:split-differential} kemudian menentukan $\delta: C \to I[1]$ sedemikian sehingga
% segment-id: o014.li2.chapter3.resolutions.g013
	\begin{equation}\label{eqn:horseshoe-aux}\begin{tikzcd}
		L^n \arrow[r, "{\Phi^n}", "\sim"'] \arrow[d, "{d_L^n}"'] & I^n \oplus C^n \arrow[d, "{e^n}"] \\
		L^{n+1} \arrow[r, "{\Phi^{n+1}}"', "\sim"] & I^{n+1} \oplus C^{n+1}
		\end{tikzcd} \quad \text{selalu komutatif, dengan}\;
		e^n := \begin{pmatrix} d_I^n & \delta^n \\ 0 & d_C^n \end{pmatrix}.
	\end{equation}

	Demikian pula, dengan menggunakan Lema \ref{prop:split-differential}, dari sembarang $\theta: K \to I[1]$ dapat dibangun kompleks $J$ sedemikian sehingga
% segment-id: o014.li2.chapter3.resolutions.q018
	\[ J^n := I^n \oplus K^n, \quad d_J^n := \begin{pmatrix}
		d_I^n & \theta^n \\
		0 & d_K^n
	\end{pmatrix}, \quad 0 \to I \to J \to K \to 0 \;\text{eksak}. \]
	Kita ingin memilih $\theta$ sedemikian sehingga komposisi $L^n \xrightarrow[\sim]{\Phi^n} I^n \oplus C^n \xrightarrow{(\identity, \eta^n)} J^n$ memberikan morfisme $L \to J$. Jika sifat ini terpenuhi, maka $L \to J$ merupakan monomorfisme. Tuliskan $\kappa$ untuk komposisi $B \to L \to J$; morfisme ini tetap merupakan monomorfisme. Mudah diperiksa bahwa diagram dalam pernyataan komutatif. Karena $\epsilon$ dan $\eta$ keduanya kuasi-isomorfisme, funktorialitas barisan eksak panjang dalam Proposisi \ref{prop:long-exact-sequence-ses}, bersama dengan Lema Lima (Proposisi \ref{prop:5-lemma}), menunjukkan bahwa $\kappa$ juga merupakan kuasi-isomorfisme, sehingga merupakan resolusi injektif yang dicari.

	Bagaimana memilih $\theta$? Agar $L^n \to J^n$ di atas menjadi morfisme kompleks, berdasarkan \eqref{eqn:horseshoe-aux}, syarat perlu dan cukupnya ialah
% segment-id: o014.li2.chapter3.resolutions.q019
	\[\begin{pmatrix}
		\identity_{I^{n+1}} & 0 \\
		0 & \eta^{n+1}
	\end{pmatrix} \begin{pmatrix}
		d_I^n & \delta^n \\
		0 & d_C^n
	\end{pmatrix} = \begin{pmatrix}
		d_I^n & \theta^n \\
		0 & d_K^n \\
	\end{pmatrix} \begin{pmatrix}
		\identity_{I^n} & 0 \\
		0 & \eta^n
	\end{pmatrix}, \quad n \in \Z ; \]
\footnote{Catatan penerjemah (O014-C058): entri kanan bawah matriks kiri
pada sumber tercetak sebagai $d_K^n$. Agar komposisi bertipe benar pada
$I^n\oplus C^n$ dan menghasilkan persamaan di atas, target memulihkannya
menjadi $d_C^n$.}
	dengan kata lain, $\delta = \theta \eta$. Karena kuasi-isomorfisme $\eta$ merupakan monomorfisme di $\cate{C}(\mathcal{A})$, sedangkan $I[1]$ terdiri atas objek-objek injektif, Lema \ref{prop:ext-alpha-beta} memastikan keberadaan $\theta: K \to I[1]$ semacam itu. Terbukti.
\end{bukti}

% segment-id: o014.li2.chapter3.resolutions.p010
Untuk menyatakan hasil berikutnya, bagi kompleks $X$ yang diberikan dan setiap $p \in \Z$, definisikan
% segment-id: o014.li2.chapter3.resolutions.q020
\[ B^p := \Image\left( d_X^{p-1} \right), \quad Z^p := \Ker\left( d_X^p \right), \quad H^p := Z^p/B^p = \Hm^p(X) , \]
sehingga $X$ terurai menjadi barisan-barisan eksak pendek berikut:
% segment-id: o014.li2.chapter3.resolutions.q021
\[ 0 \to Z^p \to X^p \xrightarrow{d_X^p} B^{p+1} \to 0, \quad 0 \to B^p \to Z^p \to H^p \to 0. \]
Resolusi Cartan--Eilenberg yang akan diperkenalkan dapat dibayangkan sebagai resolusi serempak bagi barisan-barisan eksak pendek ini.

% segment-id: o014.li2.chapter3.resolutions.th003
\begin{teorema}[H.\ Cartan, S.\ Eilenberg]\label{prop:CE-resolution}
	\index{resolusi!Cartan--Eilenberg}
	Misalkan $\mathcal{A}$ memiliki cukup banyak objek injektif. Untuk setiap objek $X$ di $\cate{C}^+(\mathcal{A})$, terdapat bikompleks $I$ yang memenuhi syarat-syarat berikut, beserta morfisme $\epsilon: X \to (I^{\bullet, 0}, \dhori^{\bullet, 0})$ di $\cate{C}(\mathcal{A})$, dengan $\dhori$ dan $\dvert$ berasal dari struktur bikompleks $I$ (Definisi \ref{def:double-cplx}):
% segment-id: o014.li2.chapter3.resolutions.l007
	\begin{enumerate}[(i)]
% segment-id: o014.li2.chapter3.resolutions.i013
		\item Untuk semua $(p, q) \in \Z^2$, berlaku $q < 0 \implies I^{p, q} = 0$.
% segment-id: o014.li2.chapter3.resolutions.i014
		\item Ambil $N \in \Z$ sedemikian sehingga $n < N \implies X^n = 0$. Maka, untuk semua $(p, q) \in \Z^2$, berlaku $p < N \implies I^{p, q} = 0$.
% segment-id: o014.li2.chapter3.resolutions.i015
		\item Untuk setiap $p \in \Z$, kita memiliki resolusi injektif dari $X^p \in \Obj(\mathcal{A})$,
% segment-id: o014.li2.chapter3.resolutions.q022
		\[ 0 \to X^p \xrightarrow{\epsilon^p} I^{p, 0} \xrightarrow{\dvert^{p, 0}} I^{p, 1} \xrightarrow{\dvert^{(p, 1)}} \cdots . \]
% segment-id: o014.li2.chapter3.resolutions.i016
		\item Morfisme di atas menginduksi resolusi injektif dari $Z^p := \Ker\left( d_X^p \right)$,
% segment-id: o014.li2.chapter3.resolutions.q023
		\[ 0 \to \Ker\left( d_X^p \right) \to \Ker\left( \dhori^{p, 0} \right) \to \Ker\left( \dhori^{p, 1} \right) \to \cdots . \]
% segment-id: o014.li2.chapter3.resolutions.i017
		\item Demikian pula, $B^{p+1} := \Image\left( d_X^p \right)$ memiliki resolusi injektif
% segment-id: o014.li2.chapter3.resolutions.q024
		\[ 0 \to \Image\left( d_X^p \right) \to \Image\left( \dhori^{p, 0} \right) \to \Image\left( \dhori^{p, 1} \right) \to \cdots . \]
% segment-id: o014.li2.chapter3.resolutions.i018
		\item Demikian pula, $H^p := \Hm^p\left(X\right)$ memiliki resolusi injektif
% segment-id: o014.li2.chapter3.resolutions.q025
		\[ 0 \to \Hm^p\left(X\right) \to \underbracket{\Hm^p\left( I^{\bullet, 0}, \dhori \right) \to \Hm^p\left( I^{\bullet, 1}, \dhori \right) \to \cdots}_{= \Hm^p_{\mathrm{I}}\left( I \right), \;\text{lihat \S\ref{sec:double-cplx-coh}} } . \]
	\end{enumerate}
	Data $(I, \epsilon)$ dengan sifat-sifat di atas disebut \emph{resolusi Cartan--Eilenberg} dari $X$.
\end{teorema}
% segment-id: o014.li2.chapter3.resolutions.pf007
\begin{bukti}
	Ambil $N \in \Z$ sedemikian sehingga $n < N \implies X^n = 0$. Sebelumnya telah didefinisikan barisan-barisan eksak pendek
% segment-id: o014.li2.chapter3.resolutions.q026
	\begin{align*}
		0 \to Z^N \to X^N \to B^{N+1} \to 0, & & 0 \to B^{N+1} \to Z^{N+1} \to H^{N+1} \to 0, \\
		0 \to Z^{N+1} \to X^{N+1} \to B^{N+2} \to 0, & & 0 \to B^{N+2} \to Z^{N+2} \to H^{N+2} \to 0, \\
		\vdots & &
	\end{align*}
	Untuk setiap $H^n$ (atau $B^n$), pilih resolusi injektif dengan cara seperti dalam Contoh \ref{eg:resolution-single}; untuk $n < N$ (atau $n \leq N$), ambillah resolusi tersebut sebagai $0 \to 0$.
% segment-id: o014.li2.chapter3.resolutions.l008
	\begin{compactitem}
% segment-id: o014.li2.chapter3.resolutions.i019
		\item Perhatikan bahwa $Z^N = H^N$. Proposisi \ref{prop:horseshoe} memperluas resolusi injektif dari $Z^N$ dan $B^{N+1}$ secara bersamaan menjadi resolusi injektif $I^{N, \bullet}$ dari $X^N$, yang kompatibel dengan barisan eksak pendek pertama.
% segment-id: o014.li2.chapter3.resolutions.i020
		\item Selanjutnya, gunakan Proposisi \ref{prop:horseshoe} untuk memperluas resolusi injektif dari $B^{N+1}$ bersama dengan resolusi injektif dari $H^{N+1}$ menjadi resolusi injektif dari $Z^{N+1}$, yang kompatibel dengan barisan eksak pendek kedua.
% segment-id: o014.li2.chapter3.resolutions.i021
		\item Sekarang ulangi operasi yang sama untuk memperoleh resolusi injektif $I^{N+1, \bullet}$ dari $X^{N+1}$ beserta resolusi injektif dari $Z^{N+2}$, yang masing-masing kompatibel dengan barisan eksak pendek ketiga dan keempat.
	\end{compactitem}

	Dengan melanjutkan cara ini, untuk setiap $n$ diperoleh resolusi injektif $X^n \to I^{n, 0} \to I^{n,1} \to \cdots$; untuk $n < N$, tetapkan $I^{n, \bullet} := 0$. Konstruksi di atas juga menentukan cara mendefinisikan $I^{n,m} \to I^{n+1,m}$ untuk memperoleh bikompleks $(I, \dhori, \dvert)$; perlu diperiksa bahwa $\dhori^2 = 0$, tetapi hal ini tidak mengandung kesulitan hakiki.
\footnote{Catatan penerjemah (O014-C059): sumber mengulang pemetaan vertikal
$I^{n,m}\to I^{n,m+1}$ di tempat konstruksi diferensial horizontal.
Pemeriksaan $\dhori^2=0$ pada kalimat berikutnya menetapkan arah yang dimaksud,
yakni $I^{n,m}\to I^{n+1,m}$, yang dipulihkan dalam target.}
\end{bukti}

% segment-id: o014.li2.chapter3.resolutions.p011
$\epsilon$ juga dapat dipandang sebagai morfisme bikompleks $X \to I$, asalkan $X$ diidentifikasi dengan bikompleks yang terkonsentrasi pada baris ke-$0$ (yakni sumbu horizontal $q=0$).

% segment-id: o014.li2.chapter3.resolutions.rm001
\begin{catatan}\label{rem:CE-split}
	Untuk semua $p, q \in \Z$, barisan-barisan eksak pendek yang terlibat dalam resolusi Cartan--Eilenberg,
% segment-id: o014.li2.chapter3.resolutions.q027
	\begin{gather*}
		0 \to \Ker\left( \dhori^{p,q} \right) \to I^{p,q} \xrightarrow{\dhori^{p,q}} \Image\left( \dhori^{p,q} \right) \to 0, \\
		0 \to \Image\left( \dhori^{p-1,q} \right) \to \Ker\left( \dhori^{p,q} \right) \to \underbracket{\Hm^p\left( I^{\bullet, q}, \dhori \right)}_{=: \Hm_{\mathrm{I}}(I)^{p,q}} \to 0, \\
		0 \to \Image\left( \dhori^{p-1,q} \right) \to I^{p, q} \to \Coker\left( \dhori^{p-1,q} \right) \to 0,
	\end{gather*}
	semuanya memiliki objek injektif pada suku paling kiri. Oleh sebab itu, Lema \ref{prop:inj-proj-split} mengakibatkan bahwa barisan-barisan tersebut terbelah. Ingat bahwa barisan eksak pendek terbelah dipertahankan oleh setiap funktor aditif $F: \mathcal{A} \to \mathcal{B}$. Untuk $q$ tetap, citra barisan-barisan eksak pendek di atas di bawah $F$ memberikan penguraian yang bersesuaian bagi kompleks $(F(I^{\bullet, q}), F(\dhori))$. Secara khusus, dari sini tampak bahwa
% segment-id: o014.li2.chapter3.resolutions.q028
	\[ F\Hm^p\left( I^{\bullet, q}, \dhori \right) \simeq \Hm^p\left( F(I^{\bullet, q}), F(\dhori) \right). \]
	Dengan kata lain, $F$ mempertahankan kohomologi pada arah horizontal dalam resolusi Cartan--Eilenberg.
\end{catatan}

% segment-id: o014.li2.chapter3.resolutions.rm002
\begin{catatan}\label{rem:double-cplx-resolution}
	Resolusi Cartan--Eilenberg memberikan cara lain yang tidak langsung untuk membangun resolusi injektif bagi $X \in \Obj(\cate{C}^+(\mathcal{A}))$. Pertama, pandang $X$ sebagai bikompleks yang terkonsentrasi pada baris ke-$0$. Mudah dilihat bahwa $\epsilon: X \to I$ menjadi morfisme di $\cate{C}^2_f(\mathcal{A})$. Dalam notasi \S\ref{sec:double-cplx-coh}, ambil terlebih dahulu kohomologi pada arah vertikal $\Hm_{\mathrm{II}}$ dari $X$ dan $I$, lalu kohomologi pada arah horizontal $\Hm_{\mathrm{I}}$. Morfisme terinduksi $\Hm_{\mathrm{I}} \Hm_{\mathrm{II}} (X) \to \Hm_{\mathrm{I}} \Hm_{\mathrm{II}} (I)$ merupakan isomorfisme. Dengan demikian, Teorema \ref{prop:double-cplx-tot} menunjukkan bahwa
% segment-id: o014.li2.chapter3.resolutions.q029
	\[ \tot(\epsilon): X = \tot(X) \to \tot(I) \]
	merupakan kuasi-isomorfisme di $\cate{C}^+(\mathcal{A})$. Perhatikan bahwa $\tot^n(I)$ merupakan jumlah langsung hingga dari objek-objek injektif untuk setiap $n \in \Z$. Dengan demikian, kita memperoleh resolusi injektif $X \to \tot(I)$.
\end{catatan}
	
